\section{System Modelling}\label{sec:model}\label{sec:intro}
\subsection{Scope and Assumptions}
This report identifies the Kitticopter altitude plant and evaluates its three-op-amp proportional controller. The DAQ interface connects the physical circuit to a simulated plant. The required outcomes are tracking accuracy above 92\%, overshoot below 30\%, and robustness to $\pm10\%$ aerodynamic and component uncertainty \cite{lab}.

The laboratory model uses unit mass, constant gravity and drag proportional to vertical velocity. Upward altitude is positive; thrust is linearised around hover. The experimentally unknown quantities are velocity gain $A$, time constant $\tau$ and sensor coefficient $k_s$.
\subsection{Forces and Differential Equation}
\Fref{fig:fbd} defines the force directions. For upward motion, thrust opposes gravity and drag; drag reverses when velocity reverses.
\begin{figure}[H]\raggedright
\href{\RepoRoot/tree/\RepoBranch/report/figures}{\color{black}\begin{tikzpicture}[>=Latex,scale=.66]
\draw[fill=AngelBlue!28,rounded corners] (-1,-.3) rectangle (1,.3);
\node at (0,0){$m=1$};
\draw[->,thick] (0,.3)--(0,1.55) node[above]{$f_t(t)$};
\draw[->,thick] (-.55,-.3)--(-.55,-1.5) node[below left]{$g=9.81\unit{m/s^2}$};
\draw[->,thick] (.75,-.3)--(.75,-1.5) node[below right]{$b\dot y(t)$};
\draw[->] (2,-.8)--(2,1) node[above]{$+y$};
\end{tikzpicture}}
\caption{Forces on the mass-normalised Kitticopter.}\label{fig:fbd}
\end{figure}
Newton's second law and the local thrust relation give
\begin{align}
\ddot y&=f_t-b\dot y-g+d_a,\qquad f_t=g+k_T(u-u_h),\label{eq:forces}\\
\ddot y+b\dot y&=k_T\delta u+d_a,\qquad \delta u=u-u_h,\quad u_h=2.5\unit{V}.\label{eq:dynamic}
\end{align}
Here $y$ is altitude [m], $u$ is command [V], $b$ is drag [s$^{-1}$], $k_T$ is thrust slope [m/(s$^2$ V)] and $d_a$ is acceleration disturbance [m/s$^2$]. Gravity cancels at the correctly calibrated hover point.
\subsection{Transfer Functions}
With zero initial conditions and no disturbance, taking Laplace transforms yields
\begin{align}
\frac{V(s)}{\Delta F_t(s)}&=\frac1{s+b},\qquad
G_v(s)=\frac{V(s)}{\Delta U(s)}=\frac{k_T}{s+b}=\frac{A}{\tau s+1},\label{eq:velocityplant}\\
G_y(s)&=\frac{Y(s)}{\Delta U(s)}=\frac{A}{s(\tau s+1)},\qquad \tau=1/b,\quad A=k_T/b.\label{eq:plant}
\end{align}
The position integrator explains why a constant command above hover produces continually increasing altitude rather than a settled position. Differentiating altitude therefore exposes a first-order response suitable for identification.
\page

\subsection{Complete Signal Architecture}
\Fref{fig:block} distinguishes the position reference from the voltage command. Both reference and altitude are scaled by $k_s$ [V/m]; the controller acts on their voltage difference.
\begin{figure}[H]\raggedright
\href{\RepoRoot/tree/\RepoBranch/report/figures}{\color{black}\begin{tikzpicture}[>=Latex,line width=.8pt,font=\small,
 block/.style={draw,fill=AngelBlue!22,minimum height=10mm,minimum width=16mm,align=left},
 sum/.style={draw,circle,inner sep=1pt,minimum size=7mm}]
\node (x) at (0,0){$x$ [m]};
\node[block] (ksr) at (1.5,0){$k_s$};
\node[sum] (err) at (3.2,0){$\Sigma$};
\node[block] (kp) at (5.1,0){$K_p$};
\node[sum] (bias) at (7.1,0){$\Sigma$};
\node[block] (sat) at (9.6,0){0--5 V};
\coordinate (u) at (12.5,0);
\draw[->] (x)--(ksr);
\draw[->] (ksr)--node[above=5.5mm]{$r_v$ [V]}node[pos=.92,below]{$+$}(err);
\draw[->] (err)--node[above=5.5mm]{$e_v$ [V]}(kp);
\draw[->] (kp)--node[above=5.5mm]{$K_pe_v$}(bias);
\draw[->] (bias)--node[above=5.5mm]{$u_c$ [V]}(sat);
\draw[->] (sat)--node[above=5.5mm]{command $u$ [V]}(u);
\draw[->] (7.1,0.95)node[above]{$u_h=2.5$ V}--node[right]{$+$}(bias);
\node[sum] (sub) at (12.5,-3.0){$\Sigma$};
\node[block] (thrust) at (10.4,-3.0){$k_T$};
\node[sum] (force) at (8.2,-3.0){$\Sigma$};
\node[block] (mech) at (5.9,-3.0){$\dfrac{1}{s+b}$};
\node[block] (int) at (2.7,-3.0){$\dfrac{1}{s}$};
\coordinate (ytap) at (1.0,-3.0);
\node (y) at (-.1,-3.0){$y$ [m]};
\draw[->] (u)--node[right]{plant input}(sub);
\draw[->] (14.1,-3.0)node[above]{$u_h$ [V]}--node[above]{$-$}(sub);
\draw[->] (sub)--node[above=5.5mm]{$\delta u$ [V]}(thrust);
\draw[->] (thrust)--node[above=5.5mm]{$\delta f_t$}(force);
\draw[->] (force)--node[above=5.5mm]{$a_{\mathrm{in}}$ [m/s$^2$]}(mech);
\draw[->] (mech)--node[above=5.5mm]{speed $v$ [m/s]}(int);
\draw[->] (int)--node[below=3mm]{position}(y);
\draw[->] (8.2,-1.5)node[above,align=left]{acceleration disturbance\\$d_a$ [m/s$^2$]}--node[right]{$+$}(force);
\draw[->] (8.2,-4.2)node[below,align=left]{uncompensated gravity\\$\Delta g$ [m/s$^2$]}--node[right]{$-$}(force);
\node[block] (ksf) at (1.7,-1.5){$k_s$};
\fill (ytap) circle (1.4pt);
\draw[->] (ytap)--(1,-2.3)--(1.7,-2.3)--(ksf.south);
\draw[->] (ksf)--(3.2,-1.5)--node[pos=.5,right]{$y_v$ [V]}node[pos=.88,left]{$-$}(err);
\end{tikzpicture}}
\caption{Feedback architecture with explicit inputs, speed and position outputs. The lower row flows right to left. Hover cancels nominal gravity; $\Delta g$ is the uncompensated part.}\label{fig:block}
\end{figure}
The analogue law is
\begin{equation}\label{eq:law}
u=\operatorname{sat}_{[0,5]}\{2.5+K_p(r_v-y_v)\},\quad r_v=k_sx,\quad y_v=k_sy.
\end{equation}
The command range is 0--5 V, sensor range 0--10 V and op-amp supply $\pm15$ V. Hover is removed before the deviation plant.
\subsection{Closed Loop, Error and Disturbances}
In the unsaturated region, define $q=AK_pk_s$ and $D(s)=\tau s^2+s+q$. Substituting $\Delta U=K_pk_s(X-Y)$ into \Eref{eq:plant} gives
\begin{align}
L(s)&=K_pk_sG_y(s),\qquad \frac{Y(s)}{X(s)}=\frac{L}{1+L}=\frac{q}{D(s)},\label{eq:tracking}\\
\frac{E_y(s)}{X(s)}&=\frac{\tau s^2+s}{D(s)},\qquad E_y=X-Y,\\
\frac{Y(s)}{D_a(s)}&=\frac{\tau}{D(s)},\qquad \frac{Y(s)}{\Delta G(s)}=-\frac{\tau}{D(s)}.
\end{align}
Uncompensated gravity acts downward. For $K_p>0$ the loop is stable and Type 1, with zero ideal step error. A constant actuator bias causes a position offset:
\begin{equation}\label{eq:bias}
y_\infty-x_\infty=\frac{d_0}{k_TK_pk_s}=\frac{\delta u_b}{K_pk_s}.
\end{equation}
At 5 m, $K_p=0.01432$ and an 8\% error allowance require $|\delta u_b|<5.327$ mV: hover calibration and matched scaling are essential.
\page

\section{System Identification}\label{sec:id}
\subsection{Step-Test Circuit and Recorded Input}
\Fref{fig:stepcircuit} shows the Session 1 identification circuit and measured input. The 10 k$\Omega$ potentiometer wiper feeds the DAQ ADC; the software applies its logged command to the plant. Session 2 instead uses the complete controller circuit.
\begin{figure}[H]\raggedright
\begin{minipage}[t]{.32\linewidth}\vspace{0pt}
\begin{circuitikz}[american,scale=.72,transform shape]
\draw (0,4) node[above]{$+15$ V} -- (0,3.7) to[R,l_=$R_1\ 10\,\mathrm{k}\Omega$] (0,2.4)
 to[R,l_=$P1\ 10\,\mathrm{k}\Omega$] (0,0) node[ground]{};
\draw[<-] (0,1.2)--(1.3,1.2) node[right,align=left]{ADC\\input};
\end{circuitikz}
\end{minipage}\hfill
\begin{minipage}[t]{.65\linewidth}\vspace{0pt}
\href{\RepoRoot/tree/\RepoBranch/analysis/modules}{\includegraphics[width=\linewidth]{step_input_transition.pdf}}
\end{minipage}
\caption{Identification circuit (left) and logged input transition (right). \RepoLink{data/session1/evidence}{Original circuit}; \RepoLink{analysis/modules}{Python}, \RepoLink{matlab/plots}{MATLAB}, \RepoLink{results}{data}.}\label{fig:stepcircuit}\label{fig:inputtransition}
\end{figure}
The command changes from hover near 2.5 V through 2.908 V at 15.8 s and 4.607 V at 15.9 s to 5.000 V at 16.0 s. The increment is $\Delta u=2.5$ V. Its short manual transition is retained and checked against an ideal-step approximation, since $\tau\approx8.1$ s.
\subsection{Data Preparation}
The \RepoLink{data/session1}{unmodified CSV} has 1,634 rows. A clock reset separates 31 preliminary rows from the retained 1,603 samples (0.1--160.3 s at 0.1 s spacing). This is one trial; the reattached CSV is byte-identical. The 10 V sensor clips, so velocity uses the unclipped metre channel.
\subsection{Parameter Calculations}
A backward difference removes the position integrator:
\begin{equation}\label{eq:diff}
v_i=\frac{y_i-y_{i-1}}{t_i-t_{i-1}}.
\end{equation}
For example, $(2.703-2.256)/0.1=4.47$ m/s at the 16.95 s interval midpoint. The 404 tail differences with endpoint time $t\ge120$ s average to 35.68748 m/s, giving
\begin{equation}\label{eq:gainid}
A=\frac{v_\infty-v_0}{\Delta u}=\frac{35.68748}{5.0-2.5}=14.27499\unit{m/(s\,V)}.
\end{equation}
Sensor calibration uses 22 nonzero, unsaturated pairs ($0<y_v<9.9$ V):
\begin{equation}\label{eq:sensor}
k_s=\frac{\sum_i y_i y_{v,i}}{\sum_i y_i^2}=0.9300026\unit{V/m}.
\end{equation}
The ratio $2.513/2.703=0.92971$ V/m agrees with this regression.
\page

\subsection{Time Constant, Sensor and Validation}
For a first-order step, $v_{63}=0.632v_\infty=22.55448$ m/s. Interpolating between 22.53 m/s at 24.0 s and 22.69 m/s at 24.1 s gives
\begin{align}
t_{63}&=24.0+\frac{22.55448-22.53}{22.69-22.53}(0.1)=24.01530\unit{s},\\
\tau_{63}&\approx24.01530-15.9=8.11530\unit{s}.
\end{align}
An exponential fit over 17--60 s gives $\tau=8.10393$ s and onset $t_0=15.91584$ s. Midpoint timing shifts this onset without changing $\tau$.
\plot{velocity_identification}{.64}{Differentiated altitude, adopted first-order model and model driven by the recorded finite input transition.}{fig:velocity}
Adopting the submitted values gives
\begin{equation}\label{eq:identified}
G_v(s)=\frac{14.2748}{8.1s+1},\quad G_y(s)=\frac{14.2748}{s(8.1s+1)},\quad k_s=0.93.
\end{equation}
Thus $b=1/8.1=0.1234568$ s$^{-1}$ and $k_T=14.2748/8.1=1.762321$ m/(s$^2$ V). Full substitutions and retained samples are in the \RepoLink{calculations}{worked calculation supplement}.
\subsection{Model Validation and Parameter Check}\label{sec:validation}
The rounded model is checked at $t>60$ s, after fitting 17--60 s. Sensor, residual and integrated-position plots are in the \RepoLink{report/figures}{figure folder}. The hold-out uses the same trial, so independent validation remains desirable.
\begin{table}[H]\raggedright\small
\caption{Model validation and parameter-check scores.}\label{tab:validation}
\begin{tabularx}{\linewidth}{Xl}\toprule\thead
Quantity & Result\\\midrule
Velocity RMSE, fit / held-out interval & 0.00533 / 0.00436 m/s\\
Held-out position RMSE & 1.686 m\\
Measured-input velocity RMSE, 17--60 s / $>60$ s & 0.08114 / 0.00436 m/s\\
Submitted $A=14.2748$ / $\tau=8.1$ / $k_s=0.93$ & 100\% / 96\% / 100\%\\\bottomrule
\end{tabularx}\end{table}
Integrating $\dot v=[A(u(t)-2.5)-v]/\tau$ from $v(15.7)=0$ gives transient RMSE 0.23\% of final velocity. The laboratory scores in \Fref{fig:paramcheck}, \Aref{app:evidence}, are not confidence intervals.
\page

\section{Controller Design}\label{sec:design}
\subsection{Root-Locus Gain Selection}
With sensor scaling included, the characteristic equation is
\begin{equation}\label{eq:characteristic}
1+K_pk_sG_y(s)=0\quad\Rightarrow\quad8.1s^2+s+13.275564K_p=0.
\end{equation}
The open-loop poles are 0 and $-1/8.1=-0.1234568$, with no finite zeros. The real-axis locus lies between these poles. Its asymptotes are $\pm90^\circ$, with centroid and breakaway at $s_b=-1/(2\tau)=-0.0617284$; the breakaway gain is 0.00232489.
Matching to $s^2+2\zeta\omega_ns+\omega_n^2$ gives
\begin{align}
\omega_n&=\sqrt{AK_pk_s/\tau},\quad \zeta=\frac1{2\sqrt{\tau AK_pk_s}},\quad \sigma=\frac1{2\tau},\\
M_p&=100e^{-\pi\zeta/\sqrt{1-\zeta^2}}<30\%\Rightarrow\zeta>0.357857,\label{eq:overshoot}\\
K_p&<\frac1{4\tau Ak_s\zeta_{\min}^2}=0.0181544.
\end{align}
The implemented gain 0.01432 satisfies the nominal bound. \Fref{fig:locus} checks the selected poles; \Aref{app:calculations} gives the magnitude check, and the \RepoLink{matlab/design}{MATLAB design folder} contains the root-locus method.
\plot{root_locus}{.76}{Root locus and selected gain. The complex branches have fixed real part, limiting achievable settling speed.}{fig:locus}
\begin{table}[H]\raggedright\small
\caption{Manual predictions at $K_p=0.01432$.}\label{tab:design}
\begin{tabularx}{\linewidth}{Xl}\toprule\thead
Quantity & Value\\\midrule
$\omega_n$ / $\zeta$ / $\omega_d$ & 0.153199 / 0.402930 / 0.140212\\
Closed-loop poles [s$^{-1}$] & $-0.0617284\pm j0.140212$\\
Overshoot / peak altitude for 5 m & 25.0804\% / 6.25402 m\\
Peak time $\pi/\omega_d$ / actual 2\% settling & 22.406 s / 54.89 s\\\bottomrule
\end{tabularx}\end{table}
Proportional gain cannot meet the earlier 30 s decay-envelope target because $\sigma$ is fixed. The $4/\sigma=64.8$ s approximation is conservative: actual 2\% reference settling is 54.89 s.
\page

\subsection{Original Simulink Model and Simulation Tests}\label{sec:simulation}
\Fref{fig:originalmodel} is the original supplied model. Its 0.93 gains scale reference and altitude, 0.01432 is the proportional gain, and the two 2.5 V terms add and remove hover. The velocity lag is followed by the altitude integrator.
\repfig{data/session2}{../../data/session2/simulink_model.png}{.94}{Original Simulink model. \RepoLink{matlab/simulink}{Native model-building and test scripts}.}{fig:originalmodel}
The supplied scope starts from zero altitude with a 2.5 m reference, changing to 5 m at about 2.5 s. This differs from the isolated 0-to-5 m design test. \Fref{fig:simcomparison} compares the original trace with Python using the same startup sequence. The original scope image is retained in \Fref{fig:originalscope}, \Aref{app:demo}.
\plot{simulink_comparison}{.85}{Supplied Simulink altitude and command, replotted from the scope and compared with the matching Python prediction.}{fig:simcomparison}
The supplied altitude peaks near 6.226 m (24.51\% above 5 m); the matched prediction peaks at 6.231 m. Their altitude RMSE is 0.0115 m, close to the 0.0124 m per-pixel extraction scale. The reconstructed maximum command is 2.56441 V.
For the isolated 5 m step, the manual prediction is
\begin{equation}\label{eq:exactstep}
y=5\left[1-e^{-0.0617284t}\{\cos(0.140212t)+0.44025\sin(0.140212t)\}\right].
\end{equation}
It predicts 25.08\% overshoot and command range 2.48330--2.56659 V. The reconstruction uses 0--5 V command and 0--10 V sensor limits, zero initial states and 100 s duration. Saturation dialogs and the original native model file were not supplied.
\page

\subsection{Physical Controller Schematic}\label{sec:hardware}
\Fref{fig:circuit} enlarges the KiCad stages; panels share V1 and V2. The bias/DAQ view is in \Aref{app:circuit} and the board photograph is \Fref{fig:board}.
\repfig{hardware/kicad}{kicad_signal_stages.pdf}{.97}{Enlarged KiCad stages: U1 attenuates the reference--sensor difference; U2 amplifies it; U3 adds loaded hover bias. $K_p=0.01(1.47)(0.97428)=0.0143219$. \RepoLink{hardware/kicad}{Editable full schematic and BOM}.}{fig:circuit}

\page

\section{Controller Tests}\label{sec:tests}
\subsection{Connected Test and Potentiometer Adjustment}
The complete circuit was connected: DAC-1 reference, DAC-0 sensor, ADC-1 command and ADC-2 sensor monitor, with common-ground $\pm15$ V rails. The potentiometer was adjusted until the input was approximately 2.565 V to improve tracking and minimise overshoot; its measurement node and time were not logged.
The observation is close to the predicted initial command 2.56659 V; it differs from the 2.500 V zero-error bias and terminal command 2.498 V (\Aref{app:circuit}).
\subsection{Experimental Response and Comparison}
The required tracking screenshot (\Fref{fig:demo}, \Aref{app:demo}) records the demonstration, which received full marks (20/20). At 86.10 s, reference is 5.033 m and height is 4.897 m. \Fref{fig:hardware} compares the digitised response with its model.
\plot{hardware_comparison}{.66}{Connected controller response replotted with Python and compared with the nominal prediction.}{fig:hardware}
The approximate peak is 6.229 m at 22.02 s. The terminal error and reference-based performance are
\begin{align}
e_f&=5.033-4.897=0.136\unit{m},\\
\eta&=100(1-0.136/5.033)=97.30\%,\quad M_{p,\rm ref}=100(6.229-5.033)/5.033=23.77\%.
\end{align}
\begin{table}[H]\raggedright\small
\caption{Simulation and connected-controller comparison.}\label{tab:performance}
\begin{tabularx}{\linewidth}{P{33mm}XXX}\toprule\thead
Metric & Isolated 5 m model & Original Simulink & Connected circuit\\\midrule
Reference sequence & 0 to 5 m & 2.5 to 5 m & Near 0 to 5.033 m\\
Peak altitude & 6.254 m & $\approx6.226$ m & $\approx6.229$ m\\
Peak above reference & 25.08\% & $\approx24.51\%$ & $\approx23.77\%$\\
Terminal tracking & 100\% ideal & $\approx99.98\%$ & 97.30\%\\\bottomrule
\end{tabularx}\end{table}
Terminal height passes the 8\% band but misses the 2\% band, so physical 2\% settling is unconfirmed. Overshoot normalised to terminal height is 27.20\%, also below 30\%.
\page

\subsection{Evaluation Against Specifications}\label{sec:evaluation}
\Tref{tab:evaluation} distinguishes observed nominal performance from tolerance predictions. \Fref{fig:aerodynamic} varies drag by $\pm10\%$ while holding the identified thrust slope fixed.
\begin{table}[H]\raggedright\small
\caption{Specification evaluation and limitations.}\label{tab:evaluation}
\begin{tabularx}{\linewidth}{P{34mm}Xl}\toprule\thead
Requirement & Evidence & Outcome\\\midrule
Tracking $>92\%$ & 97.30\% in the retained connected test. & Nominal pass\\
Overshoot $<30\%$ & Model 25.08\%; circuit $\approx23.77\%$ about reference. & Nominal pass\\
Drag uncertainty $\pm10\%$ & Predicted overshoot 21.15--29.45\%. & Drag-only pass\\
Components $\pm10\%$ & 2,048 corners: maximum error 48.79\%, overshoot 38.95\%. & Not guaranteed\\\bottomrule
\end{tabularx}\end{table}
\plot{aerodynamic_robustness}{.70}{Responses for nominal drag and $\pm10\%$ drag uncertainty, at unchanged thrust slope and gain.}{fig:aerodynamic}
For joint scalar drag/gain uncertainty the admissible nominal gain is
\begin{equation}\label{eq:jointbound}
K_{p,0}<\frac{(0.9b)^2}{4k_Tk_s\zeta_{\min}^2(1.1)}=0.0133683.
\end{equation}
The implemented 0.01432 exceeds this stricter bound, giving 31.42\% worst-case scalar overshoot. Independent resistor mismatch can also change the DC ratio $y_\infty/x_\infty=a_r/a_f$; nominal type number does not remove that error. \Aref{app:circuit} gives the coefficient calculations.
\subsection{Conclusion and Recommendations}\label{sec:conclusion}\label{sec:recommendations}
The controller meets nominal tracking and overshoot requirements, consistent with the full demo mark. Drag-only predictions pass; independent component tolerances and experimental repeatability remain limitations.
Use matched reference/feedback ratios, reduce both U1 1 k$\Omega$ resistors to matched 680 $\Omega$ values for $K_p\approx0.009739$, and recalibrate the loaded hover output. The corrected scalar worst-case overshoot is 23.44\%; the 5 m bias allowance becomes 3.623 mV. These changes require physical retesting. Future tests should retain repeated logs and a defined disturbance record.

